\documentclass[11pt]{article}
\usepackage[margin=1in]{geometry}
\usepackage[T1]{fontenc}
\usepackage[utf8]{inputenc}
\usepackage{lmodern}
\usepackage{amsmath,amssymb}
\usepackage[hidelinks]{hyperref}
\title{Invariance as Cause: Why Motion Changes Distance and Time}
\author{John C. W. McKinley\\Independent Researcher\\ORCID: 0009-0005-7097-5035}
\date{November 20, 2026}
\begin{document}
\maketitle
\noindent\textbf{Status:} Theoretical interpretation --- working manuscript
\begin{abstract}
Special relativity establishes that measured distances contract and elapsed proper times differ as relative velocity increases, with both effects governed by the Lorentz transformations. These relationships are mathematically well established, yet their equations describe the effects without specifying a physical purpose for their occurrence.

This paper proposes a simple governing principle: \textbf{the invariance of the speed of light, $c$, is the primary constraint, and the velocity-dependent adjustments of distance and time are the means by which the universe preserves that constraint.} Every nonzero fraction of $c$ corresponds to a definite degree of relativistic adjustment. As relative velocity increases, longitudinal distance contracts and elapsed proper travel time decreases in accordance with the Lorentz factor. These are not independent effects but coordinated expressions of a single governing requirement.

The proposal reverses the usual explanatory emphasis. Rather than treating length contraction and time dilation as consequences to be accepted through mathematical transformation alone, it interprets them teleologically as the universe's enforcement of an invariant causal speed. The equations remain unchanged; what changes is the proposed understanding of why their relationships hold.

The resulting framework introduces no modification to the predictions of special relativity. Its contribution is interpretive: a unified physical principle in which the preservation of $c$ governs the observed relationships among velocity, distance, and elapsed time.
\end{abstract}
\noindent\textbf{Keywords:} special relativity, Lorentz contraction, time dilation, invariant speed, causal constraint, teleological interpretation, Timeless Light Model

\section{Introduction}
The equations of special relativity accurately describe the relationships among velocity, distance, and elapsed time. As relative velocity increases, longitudinal distances undergo Lorentz contraction, and moving clocks exhibit time dilation when compared within a specified reference frame.

These effects are experimentally established.

Yet a distinction remains between describing a physical relationship mathematically and proposing why that relationship exists.

An equation does not govern the universe. An equation describes the behavior of the universe. The physical laws represented by the equation are what require explanation.

The present paper proposes that the invariant speed $c$ is not merely one feature of relativistic physics among several. It is the governing constraint from which the coordinated adjustments of measured distance and elapsed time should be understood.

The proposed principle is straightforward:

\begin{quote}
\textbf{The universe preserves the invariant causal speed $c$ by coordinating the measurements of distance and time across relative motion.}
\end{quote}

This is a teleological statement. It identifies the preservation of an invariant as the reason for the observed relationships, without implying that the universe possesses intentions, makes decisions, or performs calculations.

The proposal does not replace special relativity. It offers a causal-priority interpretation of its established relationships.

\section{The Invariant Constraint}
Special relativity identifies a universal invariant speed, conventionally denoted $c$.

In vacuum,
\[
c=299\,792\,458\ \mathrm{m/s}.
\]

All inertial observers obtain this same value when measuring the speed of light in vacuum, irrespective of their relative motion.

This invariance is unusual.

In ordinary Newtonian reasoning, velocities combine through simple addition. An observer moving toward an approaching object measures a greater relative speed than an observer at rest.

Light does not follow this ordinary addition rule.

A light signal measured at $c$ by one inertial observer is also measured at $c$ by another inertial observer moving relative to the first.

For that to remain true, the observers cannot retain identical assignments of spatial separation and elapsed time.

Their measurements must be related differently.

The present proposal treats this requirement as primary.

Rather than beginning with the mathematical transformations and then identifying their consequences, it begins with the invariant constraint and asks what relationships must hold for that constraint to remain satisfied.

The answer is that measurements of distance and time must change together.

\section{A Continuous Response to Relative Velocity}
Relativistic effects are not activated at some special threshold near the speed of light.

They are continuous functions of relative velocity.

Define
\[
\beta=\frac{v}{c}.
\]

For massive bodies moving at speeds below $c$,
\[
0\leq\beta<1.
\]

The Lorentz factor is
\[
\gamma=\frac{1}{\sqrt{1-\beta^2}}.
\]

The longitudinal length-contraction relationship is
\[
L=\frac{L_0}{\gamma}
=L_0\sqrt{1-\beta^2},
\]
where $L_0$ is the proper length of the measured interval.

The corresponding time-dilation relationship for a clock moving uniformly relative to an inertial frame is
\[
\Delta\tau=\frac{\Delta t}{\gamma},
\]
where $\Delta\tau$ is the elapsed proper time recorded by the moving clock between the relevant events.

Both relationships contain the same factor:
\[
\sqrt{1-\beta^2}.
\]

This shared factor is central to the proposed interpretation.

Every nonzero magnitude of relative velocity corresponds to a definite relativistic adjustment. The adjustment is extremely small at ordinary speeds and increasingly pronounced as relative velocity approaches $c$.

The proposed governing rule is therefore continuous:

\begin{quote}
\textbf{Every degree of relative motion entails the corresponding degree of adjustment in the measurements of distance and time required by the invariant causal-speed constraint.}
\end{quote}

The two adjustments are not separate interventions. They are connected expressions of one requirement.

\section{Distance and Travel Time}
Consider two stationary locations separated by a proper distance $D$ in an inertial reference frame.

A spacecraft travels between them at constant speed $v$.

In the frame where the locations are stationary, the journey takes
\[
T=\frac{D}{v}.
\]

In the spacecraft's inertial frame, the distance between the locations is contracted:
\[
D'=\frac{D}{\gamma}.
\]

The elapsed proper time aboard the spacecraft is
\[
\tau=\frac{T}{\gamma}.
\]

Substituting the original travel-time relationship gives
\[
\tau=\frac{D}{\gamma v}.
\]

Therefore,
\[
\boxed{\tau=\frac{D'}{v}}.
\]

This result demonstrates the coordinated relationship.

The spacecraft measures a shorter distance to its destination and records a correspondingly shorter travel time.

Its measured relative velocity remains
\[
\frac{D'}{\tau}=v.
\]

The velocity does not exceed $c$.

The distance and time adjustments preserve a consistent velocity relationship within the spacecraft's frame.

The proposed interpretation assigns explanatory priority to this consistency.

The universe does not permit observers to obtain contradictory values for its invariant causal speed. Distance and time measurements consequently differ between frames in precisely the relationships established by the Lorentz transformations.

\section{An Interstellar Example}
Consider a spacecraft traveling from Earth toward Alpha Centauri, approximately 4.37 light-years away.

Suppose the spacecraft maintains a speed of $0.99c$ relative to Earth.

The Lorentz factor is approximately
\[
\gamma=7.089.
\]

In Earth's reference frame, the distance is 4.37 light-years and the travel time is
\[
T=\frac{4.37}{0.99}\approx4.41\ \text{years}.
\]

In the spacecraft's frame, the distance is
\[
D'=\frac{4.37}{7.089}
\approx0.616\ \text{light-years}.
\]

The spacecraft's elapsed travel time is
\[
\tau=\frac{4.41}{7.089}
\approx0.622\ \text{years}.
\]

The spacecraft therefore reaches its destination after approximately 227 days of onboard elapsed time.

The Earth observer assigns a travel duration of approximately 4.41 years.

Neither measurement is erroneous.

The two observers assign different spatial and temporal intervals to the journey.

From the perspective of the present proposal, this is precisely the required behavior.

The invariant constraint remains satisfied because the measurements of distance and time are coordinated across the two reference frames.

The spacecraft has not exceeded the causal speed limit. It has measured a shorter distance and accumulated less proper time.

\section{Why Distance and Time Must Be Considered Together}
Length contraction alone does not establish the invariance of $c$.

Time dilation alone does not establish it either.

The complete Lorentz transformation also includes the relativity of simultaneity.

For relative motion along the $x$-axis, the transformations are
\[
x'=\gamma(x-vt)
\]
and
\[
t'=\gamma\left(t-\frac{vx}{c^2}\right).
\]

The spatial coordinate transformation depends on time, and the temporal coordinate transformation depends on position.

This interdependence is essential.

For a light signal satisfying
\[
x=ct,
\]
substitution into the Lorentz transformations gives
\[
x'=ct'.
\]

Thus,
\[
\frac{x'}{t'}=c.
\]

The invariant speed is preserved.

This is the strongest mathematical expression of the proposed governing principle.

\begin{quote}
\textbf{Space and time measurements are not independently adjustable quantities. Their transformation is coordinated so that the invariant causal-speed relationship remains unchanged.}
\end{quote}

The proposal therefore concerns the complete relationship among space, time, and relative motion, not merely the shortening of a spacecraft's journey.

\section{The Round-Trip Clock Experiment}
A further test involves a spacecraft that leaves Earth, travels to a distant destination, turns around, and returns.

At departure, the spacecraft clock and an Earth clock are synchronized.

At reunion, the spacecraft clock records less elapsed proper time.

For the idealized Alpha Centauri journey at $0.99c$, neglecting turnaround duration:

\begin{itemize}
\item Earth records approximately 8.83 years.
\item The spacecraft records approximately 1.24 years.
\end{itemize}

The difference is physical and measurable.

The two clocks began together and reunite at the same event, yet they have accumulated different elapsed proper times.

The explanation must account for this asymmetry.

During the outbound and inbound portions, the spacecraft occupies different inertial frames. The turnaround changes the spacecraft's motion and its assignment of simultaneity to distant Earth events.

The complete journey therefore cannot be represented as two observers remaining in one pair of unchanged inertial frames throughout.

The invariant-speed interpretation accommodates the result without introducing a preferred universal clock.

Each clock records the proper time associated with its own physical history.

The same relativistic rules govern both.

Their different accumulated times follow from the different histories connecting the shared departure and reunion events.

The governing constraint is universal; the resulting elapsed times are not required to be equal.

\section{The Limiting Case}
As a massive spacecraft's velocity approaches $c$,
\[
\beta\rightarrow1,
\]
and consequently,
\[
\frac{1}{\gamma}\rightarrow0.
\]

For a fixed destination distance in the destination's rest frame,
\[
D'\rightarrow0
\]
and
\[
\tau\rightarrow0.
\]

The measured distance and elapsed proper travel time approach zero together.

This limiting behavior is especially significant.

The spacecraft does not need to exceed $c$ to accumulate arbitrarily little proper travel time over a fixed journey, considered in the idealized limit.

The relevant distance contracts along with the elapsed proper time.

For a massive spacecraft, the limit is approached but not attained. Reaching $c$ would require unbounded energy under standard relativistic dynamics.

For a lightlike path, the proper-time interval is zero:
\[
d\tau=0.
\]

A photon has no valid rest frame. Consequently, the limit for massive observers must not be mistaken for an actual photon perspective.

Nevertheless, the limiting relationship reveals a direct connection between invariant causal speed, spatial separation, and elapsed proper time.

\section{The Proposed Explanatory Reversal}
The conventional mathematical presentation of special relativity begins with its postulates, establishes the Lorentz transformations, and derives the resulting relationships among measurements.

This paper proposes an explanatory reversal.

Begin with the physical requirement:
\[
\boxed{c\text{ is invariant}.}
\]

Then interpret the coordinated transformation of distance and time as the physical expression of that requirement.

The explanatory sequence becomes:

\begin{quote}
\textbf{Invariant causal-speed constraint $\rightarrow$ coordinated space-time measurement relationships $\rightarrow$ observed length contraction and time dilation.}
\end{quote}

The arrows in this sequence express explanatory priority, not a chronological process.

The universe is not assumed to detect a possible violation and subsequently correct it.

Rather, a violation is excluded by the governing physical law. The allowed relationships among measured quantities already satisfy that law.

This distinction is important.

The proposed principle is not that space first has one objectively universal length, then physically shrinks after a spacecraft accelerates.

It is that spatial and temporal measurements are necessarily dependent on the observer's state of motion.

No universal adjustment mechanism or preferred frame is introduced.

The rule is invariant; the measurements are frame-dependent.

\section{Relationship to the Timeless Light Model}
The Timeless Light Model (TLM) places particular emphasis on the causal status of $c$, rather than treating it exclusively as the velocity of a traveling photon-object.

In the TLM interpretation, $c$ represents a baseline constraint on causal registration.

The present proposal is consistent with that emphasis.

It identifies the invariance of $c$ as the governing requirement and treats the observed changes in distance and elapsed time as expressions of that requirement.

This suggests a useful conceptual separation:

\begin{enumerate}
\item \textbf{The invariant constraint:} $c$ establishes the universal causal-speed relationship.
\item \textbf{The observer-dependent measurements:} distance and elapsed time are assigned according to the applicable relativistic frame.
\item \textbf{The physical record:} clocks and measuring systems register the resulting intervals.
\end{enumerate}

This paper does not require the additional ontological commitments of TLM to establish its central interpretation.

The Lorentz relationships stand independently.

The connection to TLM lies in the explanatory priority assigned to the invariant causal constraint.

In particular, the null proper-time condition of a lightlike path is not treated as evidence of a photon experiencing an instantaneous journey.

It is treated as a limiting physical relationship that does not license an internal photon clock or photon rest frame.

The present interpretation therefore preserves the distinction between the observations made by massive observers and the null structure of lightlike relations.

\section{Scope, Consistency, and Novelty}
The proposed interpretation must satisfy several requirements.

First, it must preserve the experimentally established Lorentz transformations. It does so without alteration.

Second, it must preserve the absence of a preferred inertial reference frame. It does not introduce one.

Third, it must account for the relativity of simultaneity, rather than attributing all relativistic behavior to length contraction alone. The full transformation is retained.

Fourth, it must preserve the experimentally measured differences in elapsed proper time between clocks following different histories. The interpretation accepts those differences as physical results.

Fifth, it must distinguish between a governing constraint and an active corrective mechanism. No physical system is proposed that detects excessive velocity and modifies distances after the fact.

The paper's claim of novelty is therefore limited and precise.

The invariance of $c$, the Lorentz transformations, and the resulting length-contraction and time-dilation relationships are established physics.

The proposed contribution is to express their relationship as a single teleological principle:

\begin{quote}
\textbf{The preservation of invariant causal speed is the reason distance and time measurements must transform together.}
\end{quote}

This statement is an interpretive thesis rather than a new experimentally distinguishable dynamical law.

A stronger physical theory would require additional postulates or predictions that distinguish its mechanism from existing special relativity.

The present paper makes no such claim.

Its purpose is to provide a compact explanation of the established relationships, with invariant causal speed assigned primary explanatory status.

\section{Conclusion}
Special relativity establishes that relative motion changes the measurements of distance and elapsed time.

These changes are continuous, quantitatively linked, and governed by the Lorentz transformations.

This paper proposes a simple physical interpretation of their common origin.

\begin{quote}
\textbf{The universe preserves the invariant causal speed $c$. Distance and time measurements are necessarily coordinated so that this invariant is never violated.}
\end{quote}

Under this interpretation, length contraction and time dilation are not unrelated peculiarities of relativistic motion.

They are complementary manifestations of one governing constraint.

As relative velocity increases, longitudinal distances contract and the corresponding proper travel times decrease. In the limiting case approaching $c$, both approach zero for a fixed journey.

The equations describe these relationships.

The proposed principle explains how to think about their physical necessity.

The distinction between mathematical description and governing physical law is the central contribution of this paper.

The invariant is primary.

The transformations express it.

Distance and time conform to it.

\begin{thebibliography}{9}
\bibitem{einstein1905} Einstein, A. (1905). \emph{On the Electrodynamics of Moving Bodies.} Annalen der Physik, 17, 891--921.
\bibitem{minkowski1909} Minkowski, H. (1909). \emph{Space and Time.} Physikalische Zeitschrift, 10, 104--111.
\bibitem{taylorwheeler1992} Taylor, E. F., and Wheeler, J. A. (1992). \emph{Spacetime Physics}, 2nd ed. W. H. Freeman.
\bibitem{mckinleyminimal2026} McKinley, J. C. W. (2026). \emph{A Minimal Structural Statement of the Timeless Light Model (v3.921).} Zenodo. \url{https://doi.org/10.5281/zenodo.19167403}
\bibitem{mckinleyregistration2026} McKinley, J. C. W. (2026). \emph{c Is a Registration Bound, Not a Traveler's Speed: A Registration-Based Interpretation of the Information Speed Limit.} Zenodo. \url{https://doi.org/10.5281/zenodo.20114176}
\end{thebibliography}
\end{document}